% Translator-owned audit; not original manuscript prose.
\head{第14节中段：正则中心、纤维维数与CM平坦性}

\noindent\textbf{实际范围。}
本模块从法稿物理第85页末“Grosso modo”开始，至第88页首段“只有源为CM才可能平坦”结束；英文对应物理第55页末至第57页“More delicate”之前。法文85、86、87、88页及英文55、56、57页均实际查看页面影像；另放大法86命题14.1(b)、法87的证明和CM手插。为核对交叉引用，另实际查看法21页，不把它计入本轮翻译范围。
基线为\path{current/EGA_v21/continuations/v_dot_late.tex}的连续原段和旧注V-H-L011--028。旧TeX原字节另存于\path{inherited_fibers_exact.tex}，18条旧注另逐条保存于\path{HISTORICAL_NOTES.json}。以下所有补证及反例均为本轮译者核查，不能当成法稿原有证明。

\subsection*{V22-FIB-E01：法英对照与错误归属}

\begin{longtable}{@{}p{.15\linewidth}p{.18\linewidth}p{.56\linewidth}@{}}
\toprule
新注号 & 原页 & 本次处理\\
\midrule\endhead
FIB001 & 法85末、86首；英55--56 & 法稿voire为“甚至”，英文see为误读；维数范围交叠原样保留。\\
FIB002--004 & 法86上；英56上 & 还原$\varphi_\eta^{(m+1)}$及秩$m+1$，英文的$n+1$误写；零点在$X$上，须为中心逆像；“及平坦性”属于研究问题，不是已经得到平坦性的断言。\\
FIB005 & 法86中；英56 & 原只控制泛纤维，14.1(a)仍有特大特殊纤维反例；正文明确加强为各分支所有非空几何纤维纯$d$维。\\
FIB006 & 法86下手改；英56 & 恢复$\dim f(X)$及$p_C$的“非空”纤维；有限／拟有限跟随$f$的相应条件，不由任意$d$推出。\\
FIB007--009 & 法86末--87中；英56 & 恢复英稿遗漏的反证条件、正确纤维积和像与$L_\xi$的交；原证明的非固有闭包缺口另行补足。\\
FIB010 & 法87中；英56--57 & $f$有限／拟有限两种情形分列；拟有限且等维才推出优势，不能只用等维。\\
FIB011 & 法87；英57 & 可约CM源可能有低维像分支，另补每个分支像维数$\geq m$；反例见E05。\\
FIB012--013 & 法87下手插；英57 & 法稿“应并入14.3”及$\varphi_\xi^{(m)}$可辨；英文的问号和$C_\xi^{(m)}$不当。恢复改写为模版本的意思；模的有限性明列。\\
FIB014 & 法88首；英57 & 法稿写$X$不可约，英文误写$f$；平坦性的必要CM条件另给具体反例。\\
\bottomrule
\end{longtable}

法86(b)的修改层相当密集。可辨的有效添加包括左缘$f(X)$、“les fibres non vides des $X_i$ sur $P$ sont toutes de dim $d$”，以及$p_C$纤维下方的“non vides”。原打字层“i.e. $p_C$ est quasi-fini”和“$f$ est propre sur $k$”被划去；有限／拟有限的对应另作手改。对于不能逐字可靠分离的覆盖笔迹，本版保存源页和已辨数学意思，不伪造一份无删改的法文誊本。

法87有两处完整可辨的删除：一处曾将(a)的解释写为“至少假设$\dim f(X)>m$”，另一处曾说“$\dim f(X)=m$的情形也可由(b)处理”。它们作为改写过程记录，不插回作者现行论证。对于“第21页上方、第5节之前”的引用，当前扫描物理21页直接以第5节开始，未见该上方注；故该原引用仍未定位。本报告的独立论证不声称就是作者拟引的旧注。

\subsection*{V22-FIB-E02：一般中心的正则性与爆破识别}

\noindent\textbf{1. 任意拟凝聚模首先得到Koszul正则性。}
局部平凡化$f^*\mathcal O_P(1)$，令生成截面为$a_0,\ldots,a_r$；在适当开覆盖上可假设$a_0=1$。
先用满秩矩阵的参数空间给$m+1$个有序的一般线性式选择标架，暂把所有矩阵系数$T_{ij}$当成独立变量。相应式为
\[
\ell_i=T_{i0}+\sum_{j=1}^rT_{ij}a_j,\qquad 0\leq i\leq m.
\]
对任意$A$-模$M$，变量代换$T_{i0}\mapsto\ell_i$为$A$上多项式环的自同构；因而$\ell_0,\ldots,\ell_m$在$M[T_{ij}]$上构成逐次单射的序列，且其Koszul复形在正次数正合。对系数环局部化到泛矩阵的函数域仍保持此正合性。

满秩矩阵空间到Grassmann空间的映射为标架丛；其泛函数域是$K=k(\eta)$的忠实平坦扩域。换基只对Koszul复形作同构，故正次同调消失从标架域下降到$K$。
这给出了正文对任意拟凝聚$F$实际使用的Koszul正则性。当$F$凝聚时，在其支撑与中心零点的交上，诺特局部有限模的Koszul正则性等价于普通正则列；但这一有限性不能删去。
零点之外有一个局部坐标为单位，Koszul复形可缩，不能再要求最终商模非零。这样明确区分了泛系数的正合性、有限模普通正则列和零点之外的约定，也保留了第20版对一般非有限模问题的界限。

\noindent\textbf{2. 中心的局部坐标。}
以下$X$有限型，故局部环诺特。对$F=\mathcal O_X$应用上述结论，中心理想在其零点处由正则列
\[
a_0,\ldots,a_m
\]
生成。关联空间在目标$\mathbb P^m$的第0个标准图上由
\[
B=A[t_1,\ldots,t_m]/(a_1-t_1a_0,\ldots,a_m-t_ma_0)
\tag{V22-FIB.1}
\]
给出。记$h_j=a_j-t_ja_0$。在中心上方的局部环中，
$(a_0,h_1,\ldots,h_m)$由$(a_0,a_1,\ldots,a_m)$作可逆三角变换得到，仍为正则列；诺特局部正则列可置换，所以$h_1,\ldots,h_m$正则，且$a_0$在$B$中为非零因子。在倒置$a_0$之后，(V22-FIB.1)变为$B[1/a_0]\simeq A[1/a_0]$；$a_0$的非零因子性质使$B\to B[1/a_0]$单射，而其像正是$A[a_1/a_0,\ldots,a_m/a_0]$。因此
\[
B\simeq A[I/a_0],\qquad I=(a_0,\ldots,a_m).
\]
其他标准图相同，故关联空间正是$\operatorname{Bl}_I(X_K)$。
这是具体图表计算，不是直接假设爆破对任意非平坦基变换相容。爆破的Rees代数及仿射图的定义可核对
\href{https://stacks.math.columbia.edu/tag/01OF}{Stacks Project，Tag 01OF，第31.33.1定义及31.33.2引理}。

若$f(X)$的闭包维数$\leq m$，余维$m+1$的泛中心与该闭包不交。因此$I=\mathcal O_{X_K}$，爆破为恒等；这同时解释了正文的中心外同构及受限投影处处定义。
当$m=r$时，中心为空且Grassmann参数空间为一点，以上结论按空中心理解，并不产生无穷多个不同中心。

\subsection*{V22-FIB-E03：14.1(a)原泛纤维条件不足的反例}

取任意域$k$，一个$k$-有理点$p\in\mathbb P^3_k$，以及
\[
f:X=\operatorname{Bl}_p(\mathbb P^3_k)\longrightarrow\mathbb P^3_k.
\]
$X$为不可约光滑三维射影概形，$f$为双有理态射；它在像泛点处的纤维维数$d=0$。取$m=2$，则中心是$\mathbb P^3_K$的一个泛点所给的$K$-有理点$c$，$K=k(\eta)$，且$c\ne p$。
该中心在$p$附近之外，所以沿$f^{-1}(c)$再作的爆破，不改变例外除子
\[
E=f^{-1}(p)\simeq\mathbb P^2_k.
\]
记$\lambda=\langle c,p\rangle$为通过$c,p$的直线。它对应$Q(c)\simeq\mathbb P^2_K$的一个点。
扩展锥投影在该点的纤维就是$X_K\times_{\mathbb P^3_K}\lambda$，包含整个$E_K$，因而维数至少为2。实际上它由直线的严格变换及$E_K$组成，最大维数恰为2。
原14.1(a)却要求它等于
\[
\dim X-m=3-2=1.
\]
故原假设确有反例；这既不是英文转写错误，也不能仅以“原证明略写”消解。甚至$X$光滑且$f$固有仍不够，缺的是对特殊纤维的控制。

\subsection*{V22-FIB-E04：订正14.1及14.2的独立补证}

\noindent\textbf{1. 射影闭包与边界。}
对每个不可约分支$X_i$，置
\[
n_i=\dim X_i,\qquad Z_i=\overline{f(X_i)}_{\mathrm{red}},
\qquad e_i=\dim Z_i=n_i-d.
\]
最后一个等式来自各非空纤维维数$d$及有限型维数公式。有限型域上不可约概形在任意扩域后仍纯维原维数，所以扩域后$X_i$及$Z_i$的各不可约分支分别具有维数$n_i$及$e_i$；以下在$\kappa(\xi)$上使用的截面下界仍适用。由像的可构造性，可取$Z_i$的非空开集$U_i\subset f(X_i)$；记$B_i=Z_i\setminus U_i$，则$\dim B_i\leq e_i-1$。

泛中心$C=C_\eta$同时满足有限多个闭集的通常一般线性截面维数界：
\[
\dim(Z_{i,K}\cap C)=e_i-m-1\quad(e_i>m),\qquad
Z_{i,K}\cap C=\varnothing\quad(e_i\leq m),
\]
以及
\[
\dim(B_{i,K}\cap C)\leq e_i-m-2,
\]
其中负上界表示空集。该维数界由关联簇的维数计算得出：固定一个$e$维射影闭集，包含其点的余维$m+1$线性子空间之关联族，在Grassmann上的泛纤维维数为$e-m-1$；若右端为负则为空，而$e>m$时射影交维数下界保证非空。
这些结论对后续任意扩域仍成立。

\noindent\textbf{2. 通过中心的全部截面。}
令$\xi\in Q(C)$，换底到$\kappa(\xi)$并记相应余维$m$线性空间为$L$，于是$C\subset L$为超平面。
若$e_i>m$，$Z_i\cap L$的每个不可约分量$D$由$m$个线性方程截出，维数至少$e_i-m>0$。
若其维数大于$e_i-m$，其射影闭性保证$D\cap C$非空、维数至少$\dim D-1\geq e_i-m$，与中心的维数界矛盾。所以$Z_i\cap L$纯维$e_i-m$。
而若某个$D$包含于$B_i$，同理有
\[
\dim(D\cap C)=e_i-m-1>\dim(B_i\cap C),
\]
仍矛盾。因此每个$D$都遇到$U_i$，实际像$f(X_i)\cap L$非空且维数$e_i-m$。纤维维数公式给出
\[
\dim(X_i\times_P L)=e_i-m+d=n_i-m.
\tag{V22-FIB.2}
\]
这一步通过射影\emph{闭包及边界}完成，未把任意局部闭正维集当作必与超平面相交的射影闭集。

若$e_i\leq m$，中心避开$Z_i$。其到$Q(C)$的线性投影是固有的；任何纤维在$L$中为闭集，且不遇$C$，所以同时位于仿射空间$L\setminus C$中。该纤维既固有又仿射，因而有限。于是$X_i\times_P L$若非空，维数为$d$。

在(a)中，取$n=\max_i n_i$；共同$d$保证最大$n_i$也具有最大$e_i$。由于$\max e_i>m$，每个相关投影纤维都遇到最大维分支，贡献维数$n-m$。其余分支按(V22-FIB.2)或低像维情形贡献至多$n-m$。故每个纤维维数为$n-m$，且非空。
在(b)中所有$e_i\leq m$，故全部非空纤维维数为$d$。

\noindent\textbf{3. 有限性、拟有限性与优势性。}
当$\dim f(X)\leq m$时，取$f$的概形论像的射影闭包$Z$；泛中心与其支集不交，上述投影$Z_K\to Q(C)$有限。
若$f$有限，则它通过闭子概形$Z$的分解仍有限，复合有限；若$f$拟有限，则复合拟有限。
这证明14.1(b)最后的两种情况，无须错误地从任意$d$推出零维。

在14.2中$f$有限或拟有限，故$\dim f(X)=\dim X=m$。拟有限的$p_{C_\eta}$的像维数为$m$；目标是不可约$m$维射影空间，因此像稠密。若$f$有限，$p_{C_\eta}$的像还闭，故满射。这也说明“优势”与“满射”在两种情形中必须分别表述。

\subsection*{V22-FIB-E05：14.3的分支限制与完整CM补证}

\noindent\textbf{1. 低维分支反例。}
取$P=\mathbb P^3_k$，其中一个平面$H\simeq\mathbb P^2_k$及一个不属于$H$的$k$-点$p$，并置
\[
X=H\amalg\{p\}\hookrightarrow P,\qquad m=1.
\]
$X$是CM，$f$为闭浸入，故各非空$f$纤维纯0维。泛中心是一条既不包含于$H$又不经过$p$的直线。$H$分支的扩展投影到$Q(C)\simeq\mathbb P^1_K$具有一维纤维；但孤立点分支映到$\mathbb P^1_K$的一个点。
后者在对应正则局部环上被一个非零参数湮灭，绝非平坦。故总源CM及纤维的\emph{最大}维数相等，不足以推出平坦性。必须在每个支撑分支上排除像维数小于$m$的情形。

\noindent\textbf{2. 订正概形版本的图表与纯维数。}
现在采用14.1的加强条件，并设$X$为CM且每个$e_i=\dim f(X_i)\geq m$。
由于有限型域上CM概形局部等维，在一个足够小的开邻域内，可将遇到的分支总维数统一记作$n$。
泛中心的正则列使关联空间的标准图为(V22-FIB.1)；其中$m$个$h_j$为正则列。因此标准图是CM且纯$n$维；在闭点处，其局部环维数为$n$。中心外它与$X_K$同构。

中心本身在每个$n_i$维分支上为空或纯维$n_i-m-1$；其上方的关联纤维是$\mathbb P^m$，故这部分维数至多
\[
(n_i-m-1)+m=n_i-1.
\]
它不能成为维数$n_i$的额外分支。因而关联空间的分支正是中心外各分支的闭包。
在任意$\xi\in Q(C)$处，关联纤维为$X_{\kappa(\xi)}\times_P L_\xi$。
对每个分支，E04给出该截面维数$n_i-m$（$e_i=m$时为$d=n_i-m$）。
它在$X_i$上由$m$个方程定义，故每个非空分量的维数至少$n_i-m$；结合上界即为\emph{纯}维$n_i-m$。这里不仅计算了所有分支合起来的最大维数。

\noindent\textbf{3. 正则目标的局部参数。}
在关联空间的一个闭点$z$处，$\xi=p(z)$也为$K$上闭点。目标的正则局部环$R=\mathcal O_{Q(C),\xi}$维数$m$；取其一组正则参数$u_1,\ldots,u_m$。
源局部环$B$为CM、维数$n$，而
\[
\dim B/(u_1,\ldots,u_m)B=n-m.
\]
因此这些参数在$B$中构成正则列，商仍CM。逐次应用局部平坦判据，即得$B$对$R$平坦。有限型态射的平坦位置开，且域上有限型概形为Jacobson概形，所以在所有闭点成立的这项平坦性在全空间成立。平坦态射的CM源给出CM纤维。
这证明订正14.3。可核验的代数支点是
\href{https://stacks.math.columbia.edu/tag/00R4}{Stacks Project，Tag 00R4的奇迹平坦性}、
\href{https://stacks.math.columbia.edu/tag/00ML}{Tag 00ML的局部平坦判据}
以及\href{https://stacks.math.columbia.edu/tag/045Q}{Tag 045Q的CM态射定义及CM纤维性质}。
这些是本版补证的公开依据，不替代未定位的作者“第21页”引用。

\noindent\textbf{4. 作者提出的凝聚模版本。}
设$f$拟有限，$F$为凝聚CM模，$\operatorname{Supp}F$每个分支维数$\geq m$。
对中心应用E02的$F$-正则性，在(V22-FIB.1)的每个局部图中，$h_1,\ldots,h_m$构成$F[t_1,\ldots,t_m]$-正则列。
因为多项式扩张保持CM，商模即$\widetilde F(C)$是CM；中心外与原模相同，中心上方支撑的维数同样比相应支撑分支至少小1，故无额外支撑分支。
由$f$拟有限，$F$支撑分支的像维数等于其本身维数，E04的每个截面于是具有期望的纯维数。

在闭点局部，令$N=\widetilde F(C)_z$，其支撑维数为$n$，则
\[
\dim\operatorname{Supp}\bigl(N/(u_1,\ldots,u_m)N\bigr)=n-m.
\]
CM模的参数判据表明$u_1,\ldots,u_m$为$N$-正则列，商模CM。
以$R$的参数理想应用Tag 00ML：其Koszul分解给出
$\operatorname{Tor}_1^R(N,R/\mathfrak m_R)=0$，而纤维模对剩余域当然平坦，故$N$对$R$平坦。
再利用平坦位置开放及有限型Jacobson性质，将平坦性从闭点推广到所有点。在任意一点，目标正则局部环的一组参数因该平坦性而在相应有限模上构成正则列；总模CM，故参数商也为CM。这才给出各纤维模的CM性，从而得到相对CM性。
这里用到的有限模CM参数判据及多项式稳定性见
\href{https://stacks.math.columbia.edu/tag/00N2}{Stacks Project，Tag 00N2，第10.103.4命题及10.103.13引理}。
没有把普通拉回模的CM性与相对CM性混为一谈，亦未推广到任意非有限模。

\subsection*{V22-FIB-E06：CM条件不能换成纯粹维数条件}

原稿给的是必要条件论证，没有写出具体非CM概形。这里给一个可直接核验的例子。
令
\[
R=k[x^4,x^3y,xy^3,y^4]\subset k[x,y],
\qquad
\mathfrak m=(x^4,x^3y,xy^3,y^4).
\]
$R$为整环，且在$k[x^4,y^4]$上整且有限，因此$\dim R=2$，并且
$a=x^4,b=y^4$在$R_\mathfrak m$中为一组参数。
置$u=x^2y^2$。由于次数4的单项式生成元中恰缺$u$，有$u\notin R$；事实上$u\notin R_\mathfrak m$：若$s\notin\mathfrak m$且$su\in R$，则$s$的非零常数项使$su$的最低次齐次部分成为$u$的非零倍数，矛盾。
但是
\[
au=x^6y^2=(x^3y)^2\in R,\qquad
bu=x^2y^6=(xy^3)^2\in R.
\]
故$r=au$在$R_\mathfrak m/aR_\mathfrak m$中非零，却满足
\[
br=a(bu)\in aR_\mathfrak m.
\]
参数$b$在商上是零因子，所以这个二维局部整环不是CM。该论证在任意域扩张后保持。

令$X=\operatorname{Spec}R\hookrightarrow\mathbb A^4_k\subset\mathbb P^4_k$，取$m=2$。
泛中心投影由14.2为拟有限优势态射$X_K\to\mathbb P^2_K$。若它在顶点上方平坦，则目标正则、零维纤维为CM，由平坦局部CM传递得到源顶点局部环CM，与上面的直接计算矛盾。
这为原文“不能以维数条件代替CM”提供独立实例；没有把它误写成原作者已经举出的例子。

\subsection*{仍保留的边界}

本模块没有宣称原14.1的弱泛纤维前提成立；该前提已由E03反例否定。正文采用加强后的充分条件，条件必须随命题一起保留。
一般拟凝聚模只核验Koszul正则性，普通正则列及CM模补证明确使用有限性。原“第21页、第5节之前”的交叉引用仍未定位。源页的密集覆盖旧字未被猜补成连续法文。上述有限型域上补证不自动推广到任意基底或非诺特对象。

